% =====================================================================
%  Page de garde — conforme au modèle officiel ISIMS Sfax.
%
%  Les éléments graphiques (bande latérale, filigrane, logo) ont été extraits du PDF fourni par
%  l'institut : la couverture est imposée, elle ne se redessine pas « à peu près ».
%
%  Les valeurs à confirmer sont en rouge — voir rapport/PLAN.md, section « Reste à décider ».
% =====================================================================

% La couverture réclame environ 26 cm de hauteur utile, la mise en page du corps n'en offre que
% 24,7 : le bloc central débordait et LaTeX le renvoyait entier sur la page suivante, laissant une
% couverture vide sous l'en-tête. On lui donne sa propre géométrie plutôt que de rogner le contenu
% imposé par le modèle.
\newgeometry{top=1.2cm,bottom=1.2cm,left=3cm,right=2cm}

\begin{titlepage}
\thispagestyle{empty}
% L'interligne 1,5 et parskip du corps ajoutaient près de quatre centimètres cumulés ici. La
% couverture se compose serrée : ses espacements sont posés explicitement, un par un.
\singlespacing
\setlength{\parskip}{0pt}

% La bande décorative occupe ~4,2 cm une fois portée à la hauteur de la page ; tout le contenu est
% donc décalé au-delà, sinon le bloc institutionnel passe dessous.
\newlength{\bandeoffset}\setlength{\bandeoffset}{1.6cm}

% --- Décor : bande à gauche, filigrane à droite ---
\begin{tikzpicture}[remember picture, overlay]
  \node[anchor=north west, inner sep=0] at (current page.north west)
       {\includegraphics[height=\paperheight]{isims-bande-gauche}};
  \node[anchor=east, inner sep=0] at ($(current page.east)+(0,1.2cm)$)
       {\includegraphics[height=0.28\paperheight]{isims-filigrane-droite}};
\end{tikzpicture}

% --- Bandeau d'en-tête : institution / logo / nature du document ---
\noindent\hspace*{\bandeoffset}%
\begin{minipage}[t]{0.30\textwidth}
  \centering\footnotesize
  République Tunisienne\\
  Ministère de l'Enseignement Supérieur\\
  et de la Recherche Scientifique\\[1pt]
  \rule{2cm}{0.8pt}\\[1pt]
  Université de Sfax\\[1pt]
  \rule{2cm}{0.8pt}\\[1pt]
  Institut Supérieur d'Informatique\\
  et de Multimédia de Sfax
\end{minipage}
\hfill
% Le modèle officiel donne au logo 61,7 mm de large. À 0,20\textwidth il n'en faisait que 32 : il
% paraissait perdu entre deux blocs de texte qui le dominaient. Les colonnes latérales sont
% resserrées d'autant pour lui laisser la place, sans toucher au texte qu'elles portent.
\begin{minipage}[t]{0.34\textwidth}
  \centering
  \vspace{0.2cm}
  \includegraphics[width=\linewidth]{isims-logo}
\end{minipage}
\hfill
\begin{minipage}[t]{0.22\textwidth}
  \centering\footnotesize
  Mémoire d'ingénieur\\
  en Informatique\\[1pt]
  \rule{1.6cm}{0.8pt}\\[1pt]
  N\textsuperscript{o} d'ordre : \textcolor{red}{2026 - XXX}
\end{minipage}

\vspace{0.3cm}
\noindent\hspace*{\bandeoffset}\textcolor{isimsblue}{\rule{0.88\textwidth}{1.2pt}}

\vspace{1.1cm}

% Le bloc central est décalé de la même quantité que l'en-tête : il se centre dans la zone restante,
% pas dans la page, sinon le titre paraît déporté par rapport au filet ci-dessus.
\noindent\hspace*{\bandeoffset}%
\begin{minipage}{0.88\textwidth}
\centering

  {\fontsize{30}{34}\selectfont\bfseries MÉMOIRE}

  \vspace{1.0cm}
  {\footnotesize Présenté à}

  \vspace{0.4cm}
  {\bfseries\itshape L'Institut Supérieur d'Informatique\\ et de Multimédia de Sfax}

  \vspace{0.9cm}
  {\footnotesize En vue de la validation de}

  \vspace{0.4cm}
  {\bfseries\itshape DIPLÔME NATIONAL D'INGÉNIEUR EN INFORMATIQUE}

  \vspace{0.9cm}
  {\footnotesize Intitulé}

  \vspace{0.4cm}
  \rule{0.76\linewidth}{1.2pt}

  \vspace{0.55cm}
  {\Large\bfseries\itshape Système intelligent\\[2pt] de suivi de livraison}

  \vspace{0.55cm}
  \rule{0.76\linewidth}{1.2pt}

  \vspace{1.1cm}
  {\footnotesize Élaboré par}

  \vspace{0.4cm}
  {\large\bfseries\itshape Mohamed Aziz Hadjkacem}

  \vspace{1.2cm}
  {\footnotesize Soutenu le \textcolor{red}{XX/06/2026} devant le jury composé de}

  \vspace{0.7cm}
  % Noms à gauche, rôles alignés à droite — comme le modèle.
  \begin{minipage}[t]{0.52\linewidth}
    \raggedright\bfseries\itshape
    \textcolor{red}{M. XXX}\\[3pt]
    \textcolor{red}{M. XXX}\\[3pt]
    M. Bassem Bouaziz\\[3pt]
    \textcolor{red}{M. XXX}
  \end{minipage}%
  \begin{minipage}[t]{0.46\linewidth}
    \raggedleft
    Président\\[3pt]
    Rapporteur\\[3pt]
    Encadrant académique\\[3pt]
    Encadrant professionnel
  \end{minipage}

  \vspace{1.3cm}
  {\bfseries\small Année universitaire : 2025 / 2026}
\end{minipage}
\end{titlepage}

\restoregeometry
